% ===================================================================
%  TABEL Nr. 1 — De school. — Het lyceum. — Het klaslokaal.
%
%  Ceci n'est PAS une transcription : c'est une traduction, faite en
%  2026, en neerlandais d'aujourd'hui. D'ou trois differences avec
%  texto/io et texto/fr, et trois seulement :
%
%   * pas de \nl ni de \cc — la traduction ne suit aucune ligne
%     imprimee, et les lignes de ce fichier ne sont que du confort ;
%   * pas de \begin{VUpage} — elle n'a ni feuillet ni folio, et la
%     page de lecture ne lui en affiche pas ;
%   * le gras se pose sur le TERME seul, ponctuation dehors.
%
%  Tout le reste est identique, et doit l'etre : les cles « %%K » sont
%  celles de texto/io, macro pour macro pour l'apparat, et les renvois
%  \textsuperscript{(n)} visent les MEMES objets de la planche — ce
%  sont eux qui ouvrent les gros plans.
%
%  LA PREMIERE LANGUE NATURELLE DES DIX-SEPT, et la premiere depuis
%  l'esperanto ou la place des mots redevient un probleme. Les Pays-Bas
%  ont eu leurs cercles idistes des les annees dix ; la colonne se
%  traduit SUR L'IDO, comme les deux precedentes.
%
%  MAIS LE NEERLANDAIS RENVOIE SON VERBE A LA FIN, et pas partout :
%  seulement dans la subordonnee — « ...omdat de meester het bord
%  schoonmaakt » — et dans les temps composes, ou le participe passe
%  ferme la phrase. Le renvoi appartient a son substantif et ne bouge
%  pas ; c'est la phrase qui se refait autour, comme pour le hindi.
%
%  ET IL SOUDE SES COMPOSES EN UN SEUL MOT : « schoolbank »,
%  « leerlingen », « schoolbord », « inktpot ». Le gras se pose alors
%  sur le mot entier, non sur une moitie.
%
%  LES NOMS DE PERSONNE SONT NEERLANDISES, comme dans les quinze
%  autres colonnes : Ioannes est Jan, Iakobus est Jacob, Karolus est
%  Karel. Seuls le chien Medor et l'ane Aliboron gardent leur nom. Les
%  noms latins de la note (*) ne bougent pas : c'est d'eux que la note
%  parle.
%
%  DEUX CALQUES DU FAC-SIMILE, OTES EN 2026. Le premier est un mot :
%  trente-deux « ziehier », dont six « ziehier dat » qui rendaient le
%  « voici que » francais. « ziehier » n'est pas fautif ; il est du
%  neerlandais ecrit ancien, et on ne l'ecrit plus pour dire ou se
%  tient une chose. Le presentatif seul prend un verbe de position ou
%  de mouvement — « staat », « ligt », « zit », « hangt », « komt »,
%  « lopen » — et « ziehier dat » prend un adverbe de temps :
%  « daar gaat hij », « nu komen ze terug ».
%
%  ATTENTION AU VERBE QUE LA TOURNURE EMPORTAIT. « ziehier dat »
%  ouvrait une subordonnee, donc verbe a la fin ; l'oter sans refaire
%  la phrase laisse « want de zwaluwen zich al op de daken
%  verzamelen », qui n'est pas du neerlandais. Quatre phrases ont ete
%  reprises a ce titre.
%
%  LE SECOND CALQUE EST UN BLANC : cinquante et une fois, le
%  deux-points portait l'espace francaise qui le precede — « wij zien
%  : een meter ». Le neerlandais ne met pas cette espace, et aucune
%  des autres colonnes germaniques ou romanes ne la portait ; elle
%  n'etait passee que dans le neerlandais et le suedois, releves l'un
%  apres l'autre. Une typographie se calque aussi silencieusement
%  qu'une tournure.
%
%  ET UN TROISIEME CALQUE, TROUVE AU SECOND PASSAGE : « zeer ». Cent
%  neuf fois, contre quatre « heel » et quatre « erg ». « zeer » n'est
%  pas fautif, mais c'est l'intensif de l'ecrit soigne ; celui qu'on
%  parle est « heel », ou « erg » quand l'adjectif est desagreable —
%  « erg koud », « erg moeilijk », « heel mooi », « heel tevreden ».
%  Quatre « zeer » adverbiaux, apres le verbe, demandaient autre chose
%  qu'un mot : « boeide mij zeer » devient « boeide mij enorm »,
%  « vermaakte mij zeer » devient « vond ik heel leuk ».
%
%  ET LE FAC-SIMILE FAISAIT REGNER L'HIVER. « want de winter heerst »
%  traduit l'ido « nam regnas la vintro » — mais le francais du meme
%  bloc dit « car nous sommes en hiver », et c'est lui qui a raison
%  pour le neerlandais : « want het is winter ». L'espagnol et le
%  portugais gardent leur « reina el invierno », qui se dit ; le
%  neerlandais ne fait pas regner le temps qu'il fait.
% ===================================================================

%%K t01-apar-0 apar
\VUsaut{2.0mm}
\VUcentre{12.6pt}{120}{VERKLARENDE HANDLEIDING}
\VUsaut{3.2mm}
\VUcentre{8.4pt}{160}{BIJ}
\VUsaut{2.6mm}
\VUtitre{15.6pt}{0}{de Delmas-hulptabellen}
\VUsaut{3.4mm}
\VUornamento{9pt}
\VUsaut{5.0mm}
\VUcentre{13.2pt}{90}{EERSTE REEKS}
\VUsaut{3.0mm}
\VUfilet{16mm}
\VUsaut{5.6mm}

%%K t01-apar-1 apar
\VUtitre{15.0pt}{60}{TABEL Nr. 1}
\VUsaut{1.4mm}
\VUtitre{11.4pt}{0}{De school, het lyceum, het klaslokaal.}
\VUsaut{2.2mm}
\VUfilet{20mm}
\VUsaut{6.4mm}

%%K t01-tit-1 sub
\VUpk{10.2pt}{40}{Algemene beschrijving.}
\VUsaut{3.0mm}

%%K t01-01-1 p
1. --- Deze tabel stelt een \VUgras{klaslokaal} in een \VUgras{lyceum} voor. In dit klaslokaal
staan acht \VUgras{tafels} \textsuperscript{(18)} en acht
\VUgras{banken} \textsuperscript{(32)}. Op elke bank zitten drie leerlingen. De
\VUgras{onderwijzer} \textsuperscript{(7)} zit op een \VUgras{stoel}
\textsuperscript{(8)} achter zijn \VUgras{katheder} \textsuperscript{(6)}.

%%K t01-02-1 p
2. --- Links van de katheder staat een \VUgras{kast} \textsuperscript{(15)} met
\VUgras{glazen deuren}. Rechts is het \VUgras{schoolbord}
\textsuperscript{(1)} met de \VUgras{spons} \textsuperscript{(2)} en het
\VUgras{krijt} \textsuperscript{(3)}.

%%K t01-02-2 p
Boven de katheder hangen \VUgras{wandplaten} \textsuperscript{(9, 11, 12)} en een
\VUgras{kaart} \textsuperscript{(10)}.

%%K t01-03-1 p
3. --- Niet alle leerlingen zitten op hun \VUgras{plaats} \textsuperscript{(17)}. Jan
\textsuperscript{(4)} staat bij het bord, hij houdt de spons vast en wist de
\VUgras{meetkundige figuren} uit.

%%K t01-03-2 p
Piet staat bij de katheder, hij verzamelt de \VUgras{schriftelijke oefeningen}
\textsuperscript{(5)} en brengt ze naar de onderwijzer.

%%K t01-04-1 p
4. --- Deze onderwijzer is de leraar in de \VUgras{levende talen}. Hij leert de leerlingen
vreemde talen spreken, lezen en schrijven \VUgras{(Duits, Engels, Spaans, Italiaans, Russisch}
of \VUgras{Frans)}.

%%K t01-05-1 p
5. --- Het klaslokaal is heel ruim en heel licht. Achterin is een breed
\VUgras{raam} met twee \VUgras{bovenlichten} \textsuperscript{(78)} en de
\VUgras{glazen deur} om het te luchten en te verlichten.

%%K t01-05-2 p
Dit klaslokaal is niet koud. In het midden staat om het te verwarmen een heel goede
\VUgras{kachel} \textsuperscript{(46)} met zijn \VUgras{pijp} \textsuperscript{(45)}.

%%K t01-05-3 p
Aan de wand zijn \VUgras{kapstokhaken} \textsuperscript{(58)} bevestigd, daaraan hangen de
petten en de jassen van de scholieren.

%%K t01-tit-2 sub
\VUsaut{5.2mm}
\VUpk{10.2pt}{40}{De speelplaats}
\VUsaut{3.6mm}

%%K t01-06-1 p
6. --- De \VUgras{klok} \textsuperscript{(63)} wijst negen uur vijf aan.

%%K t01-06-2 p
De \VUgras{pauze} \textsuperscript{(76)} is zojuist afgelopen en de les begint opnieuw. De
speelruimte ligt op de \VUgras{binnenplaats} \textsuperscript{(75)}. Wij zien enkele
leerlingen die spelen. Een
\VUgras{ondermeester} \textsuperscript{(74)} houdt toezicht op hen.

%%K t01-07-1 p
7. --- Op de binnenplaats zien wij bovendien verschillende personen, namelijk de
\VUgras{inspecteur} \textsuperscript{(73)}, het schoolhoofd
\VUgras{(rector)} \textsuperscript{(71)} en de \VUgras{censor}
\textsuperscript{(72)}.

%%K t01-07-2 p
De inspecteur inspecteert de scholen, de rector leidt het lyceum, de censor houdt toezicht op
de studies.

%%K t01-07-3 p
De vierde persoon is de \VUgras{portier} \textsuperscript{(77)}. Hij draagt onder zijn arm een
register om de afwezigen in te schrijven.

%%K t01-08-1 p
8. --- Achterin op de binnenplaats staan de \VUgras{gymnastiektoestellen}
\textsuperscript{(70)} en de werkplaats voor de \VUgras{handenarbeid} \textsuperscript{(69)}.

%%K t01-08-2 p
Rechts van de tabel beginnen wij de \VUgras{lokalen} voor \VUgras{muziek} en
\VUgras{tekenen} te zien.

%%K t01-noto-1 noto
\VUnotes{143.2mm}{%
(*) Voor de \textit{doopnamen} raadt men aan ze ook volgens de Latijnse vorm over te
schrijven. Dat is het enige middel om ontelbare afwijkingen te ontkomen. Hoe zou men trouwens
kiezen tussen \textit{Giovanni, Jean,} \textit{Johann, Jan, Hans, John, Ivan,} enz.? Zullen
wij een bijzonder woordenboek voor de doopnamen maken? Laten wij uit het Latijn hun nominatief
putten en zeggen: \VUgras{Ioannes, Ioanna; Iulius, Iulia; Iakobus; Andreas; Lukas.} (Volledige
Gram.).%
}

%%K t01-tit-3 sub
\VUpk{10.2pt}{40}{Uitvoerige beschrijving.}
\VUsaut{2.4mm}

%%K t01-tit-4 sub
\VUpk{10.2pt}{40}{De goede leerlingen.}
\VUsaut{3.0mm}

%%K t01-09-1 p
9. --- De leerlingen van de eerste bank rechts van de katheder zijn
\VUgras{Hendrik} (*), \VUgras{Stefaan} en \VUgras{Karel}.

%%K t01-09-2 p
Hendrik is heel oplettend en vlijtig, hij luistert naar de onderwijzer.

Op zijn tafel heeft hij een \VUgras{schrift} \textsuperscript{(28)}, een
\VUgras{boek} \textsuperscript{(29)}, een \VUgras{woordenboek}
\textsuperscript{(30)} en een \VUgras{pennendoos} \textsuperscript{(31)}. Zijn boek heeft veel
\VUgras{bladzijden}, elk hoofdstuk bevat verscheidene
\VUgras{alinea's} en elke \VUgras{regel} veel \VUgras{woorden}.

%%K t01-09-4 p
Zijn buurman Stefaan \textsuperscript{(24)} is ijverig. Hij schrijft in zijn schrift de les
voor de volgende keer op. Hij heeft een mooi
\VUgras{handschrift}. Zijn boek is dicht. Wij zien alleen de \VUgras{omslag}
\textsuperscript{(27)}. Naast hem ligt zijn \VUgras{passer} \textsuperscript{(26)}.

%%K t01-09-5 p
Karel \textsuperscript{(23)} houdt zijn boek in de hand, hij slaat het open om zijn les te
herlezen. Hij heeft, zoals elke leerling, een \VUgras{inktpot} \textsuperscript{(25)} voor
zich. Hij is gehoorzaam.

%%K t01-10-1 p
10. --- Aan de tafel ernaast zitten \VUgras{Lodewijk, Anton} en \VUgras{Paul}. Lodewijk zegt
zijn \VUgras{les} \textsuperscript{(22)} op en beantwoordt de
\VUgras{vragen} van de onderwijzer. Anton \textsuperscript{(21)} is heel
onoplettend, soms straft de onderwijzer hem zelfs. Paul \textsuperscript{(20)} is een goede
kameraad, vriendelijk en behulpzaam. Hij is heel oplettend.

%%K t01-tit-5 sub
\VUsaut{4.2mm}
\VUpk{10.2pt}{40}{De slechte leerlingen.}
\VUsaut{3.2mm}

%%K t01-11-1 p
11. --- Achter Hendrik zit \VUgras{Joris} \textsuperscript{(38)}. Hij is geen slechte jongen,
maar hij is \VUgras{praatziek}, onoplettend en baldadig. Zijn
\VUgras{lichtzinnigheid} en \VUgras{ongehoorzaamheid} veroorzaken
\VUgras{wanorde} tijdens de les en dikwijls verdient hij een strenge
\VUgras{waarschuwing}. Zijn \VUgras{kladschrift} \textsuperscript{(35)} is
vuil, hij \VUgras{bevlekt} het \VUgras{met inkt} met zijn \VUgras{pen} \textsuperscript{(34)},
daarom gebruikt hij altijd zijn \VUgras{vlakgom} \textsuperscript{(36)}; zijn
\VUgras{schriftenmap} \textsuperscript{(33)} is gescheurd, want hij is erg slordig. Waarom
brengt hij zijn
\VUgras{potloodhouder} \textsuperscript{(37)} mee naar het lokaal van de
levende talen?

%%K t01-11-2 p
Zijn buurman Edmund \textsuperscript{(39)} houdt niet van hem. Hij is weliswaar een beetje
traag, maar hij is zwijgzaam en rustig. Hij praat niet; maar in plaats van te luisteren leest
hij liever de \VUgras{verhalen} uit zijn
\VUgras{leesboek}.

%%K t01-11-3 p
De derde leerling van deze bank is \VUgras{Lodewijk} \textsuperscript{(40)}. Hij is vlijtig.
Hij schrijft op een wit \VUgras{blad papier}. Voor zich heeft hij zijn \VUgras{vloeipapier}
\textsuperscript{(41)} gelegd.

%%K t01-12-1 p
12. --- De leerlingen van de tweede bank, links van de onderwijzer, zijn heel jong. De eerste,
de kleine \VUgras{Emiel}, kan nog niet zo goed schrijven, hij brengt zijn \VUgras{lei}
\textsuperscript{(42)}, zijn \VUgras{griffel} en zijn
\VUgras{sponsje} \textsuperscript{(43)} mee.

%%K t01-12-2 p
Naast hem zitten \VUgras{August} en \VUgras{Bernard} \textsuperscript{(44)}. Deze laatste
werkt goed, maar zijn \VUgras{kleding} is erg verwaarloosd.

%%K t01-13-1 p
13. --- Achter hen zitten drie \VUgras{luiaards}. \VUgras{Albert} leert zijn lessen niet.
\VUgras{Jacob} \textsuperscript{(47)} slaapt in plaats van te luisteren. Zijn \VUgras{luiheid}
levert hem \VUgras{verwijten} en zelfs
\VUgras{straffen} op. Ten slotte is \VUgras{Christiaan}
\textsuperscript{(48)} te weinig ernstig. Hij \VUgras{gedraagt zich} slecht en doet geen
enkele moeite om zich te beteren.

%%K t01-14-1 p
14. --- Zijn buren daarentegen gedragen zich goed. \VUgras{Ernst} \textsuperscript{(51)} houdt
van orde en regelmaat, hij gebruikt altijd zijn
\VUgras{liniaal} \textsuperscript{(50)} en zijn \VUgras{potlood}
\textsuperscript{(49)}. Hij is een \VUgras{externe leerling}.

%%K t01-14-2 p
\VUgras{Frans} \textsuperscript{(52)} is een \VUgras{interne leerling}, hij
draagt een uniform, eet en slaapt in het lyceum. Hij is een goede
\VUgras{kameraad}.

%%K t01-14-3 p
Maurits slijpt zijn \VUgras{potlood} met zijn \VUgras{zakmesje} \textsuperscript{(56)}, hij is
heel zorgvuldig.

%%K t01-14-4 p
Hij brengt al zijn boeken mee, zijn \VUgras{spraakkunst} \textsuperscript{(55)}, zijn
\VUgras{notitieboekje} \textsuperscript{(54)}, en zelfs een \VUgras{schrijfonderlegger}
\textsuperscript{(53)}. Zijn
\VUgras{schooltas} \textsuperscript{(57)} staat op de grond achter hem.

%%K t01-14-5 p
De twee tafels achterin het klaslokaal worden door \VUgras{goede leerlingen}
\textsuperscript{(19)} bezet.

%%K t01-tit-6 sub
\VUsaut{4.6mm}
\VUpk{10.2pt}{40}{Tekenen en muziek.}
\VUsaut{3.4mm}

%%K t01-15-1 p
15. --- In het \VUgras{tekenlokaal} hangen sierlijke \VUgras{modellen} aan de muur en een
\VUgras{lamp} \textsuperscript{(62)} met \VUgras{kap}, aan het plafond opgehangen. De
\VUgras{leraar} \textsuperscript{(60)} is heel geduldig, hij verbetert de \VUgras{tekeningen}
\textsuperscript{(61)} van de leerlingen en leert hun een \VUgras{buste}
\textsuperscript{(59)} naar de natuur tekenen.

%%K t01-16-1 p
16. --- Het \VUgras{muzieklokaal} lijkt heel interessant. Daar leert men
\VUgras{muziek} lezen, de \VUgras{toonladder} \textsuperscript{(67)} solfegen
die op de \VUgras{notenbalk} geschreven staat (do, re, mi, fa, sol, la, si\,=\,C. D. E. F. G.
A. B.), de \VUgras{sleutels} kennen en bekoorlijke
\VUgras{melodieën} zingen. Men legt de muziekstukken op de \VUgras{lessenaar}
\textsuperscript{(65)}. Een van de leerlingen \VUgras{speelt piano} \textsuperscript{(64)} en
de \VUgras{muziekleraar} \textsuperscript{(66)}
\VUgras{slaat de maat} \textsuperscript{(68)}.

%%K t01-17-1 p
17. --- Wij beginnen een hoek van de \VUgras{werkplaats} voor
\VUgras{handenarbeid} \textsuperscript{(69)} te zien, waar de jongens hout
mogen leren schaven en ijzer vijlen. Dat bestaat niet in alle lycea.

%%K t01-tit-7 sub
\VUsaut{5.0mm}
\VUpk{10.2pt}{40}{Het lokaal voor de wetenschappen.}
\VUsaut{3.6mm}

%%K t01-18-1 p
18. --- De leraar wiskunde onderwijst de leerlingen
\VUgras{meetkunde, rekenkunde, het rekenen} en het \VUgras{metriek stelsel}. Wij zien op de
tabel de voornaamste figuren van de meetkunde: een \VUgras{cirkel} \textsuperscript{(a)}, een
\VUgras{vierkant} \textsuperscript{(b)}, een
\VUgras{driehoek} \textsuperscript{(c)}, een \VUgras{lijn}
\textsuperscript{(d)}.

%%K t01-18-2 p
Wij zien ook een \VUgras{berekening}; het is noch een \VUgras{optelling}, noch een
\VUgras{aftrekking}, noch een \VUgras{vermenigvuldiging}: het is een
\VUgras{deling} \textsuperscript{(e)}.

%%K t01-19-1 p
19. --- Op de tabel van het metriek stelsel zien wij: een \VUgras{meter}
\textsuperscript{(a)}, een \VUgras{weegschaal} \textsuperscript{(b)} en haar
\VUgras{gewichten}, een \VUgras{liter} \textsuperscript{(e)}, een
\VUgras{decaliter} \textsuperscript{(f)}, een \VUgras{deciliter} en een
\VUgras{kubieke meter} \textsuperscript{(d)}.

%%K t01-19-2 p
De leerlingen bestuderen ook de \VUgras{natuurwetenschappen} \textsuperscript{(11)} --- de
dierkunde \textsuperscript{(a)}, de plantkunde \textsuperscript{(b)}, de aardkunde
\textsuperscript{(c)} --- en de
\VUgras{geschiedenis} \textsuperscript{(12)}.

%%K t01-19-3 p
In de kast staan de toestellen voor \VUgras{natuurkunde} \textsuperscript{(15)} en voor
\VUgras{scheikunde} \textsuperscript{(16)}.

%%K t01-19-4 p
Boven op de kast staan: een \VUgras{bol} \textsuperscript{(14)}, een
\VUgras{kegel} en een \VUgras{wereldbol} \textsuperscript{(13)}.

%%K t01-19-5 p
Boven de tabellen hangt een \VUgras{kaart} die de vijf werelddelen toont:
\VUgras{Europa} \textsuperscript{(g)}, \VUgras{Azië} \textsuperscript{(e)},
\VUgras{Afrika} \textsuperscript{(f)}, \VUgras{Amerika}
\textsuperscript{(ab)} en \VUgras{Oceanië} \textsuperscript{(h)}, en de twee grote oceanen, de
\VUgras{Stille} \textsuperscript{(c)} en de
\VUgras{Atlantische} \textsuperscript{(d)}.
\VUsaut{6.0mm}
\VUfilet{22mm}
